Los resultados experimentales permiten concluir que la programación dinámica es la mejor alternativa para resolver el problema AniMarathon en la práctica, ya que garantiza soluciones óptimas con un costo computacional razonable para tamaños moderados. 

La fuerza bruta, si bien entrega resultados óptimos, resulta inviable para instancias grandes debido a su complejidad exponencial. Por otro lado, los algoritmos greedy destacan por su eficiencia, pero presentan pérdidas significativas en calidad de solución en múltiples casos.

En consecuencia, la elección del algoritmo depende del contexto: si se requiere exactitud, se recomienda programación dinámica; si se prioriza rapidez, los métodos greedy son adecuados.